\chapter{Introduction}\label{sec:1}

Meaning variation and change of words or complex linguistic expressions are an integral part of language (change) (\citealt{Blank1997}, \citealt{Traugott1989}, \citealt{Eckardt2006}, among others). However, the question as to how meaning change and variation exactly work is still an unresolved issue in linguistic research. Describing the meaning change of words in time and space and the mechanism underlying it is a challenging task as the historical and/or the cognitive or psychological conditions under which meaning change happens are not always given or known. This is probably one of the reasons why there are different views on meaning change and variation. Some linguists take a more algorithmic and logical perspective on describing this mechanism, assuming that meaning change is based on a compositional reanalysis of the prior meaning into a new meaning, that is, the reanalysis can be made fully transparent and explainable by formal semantic principles as argued by formal diachronic semanticists (\citealt{Eckardt2006}, \citealt{Bar-asher-siegal2020}, \citealt{Condoravdi2015}, \citealt{Aloni2021}, among many others). According to these authors, it is possible to describe the reanalysis involved in meaning change on the level of morphosyntactic form (F) and the level of meaning (M) (see also \citealt{DetgesEtAl2021}). The important claim of formal diachronic semantics on meaning change is that at both points of time (t\textsubscript{1} and t\textsubscript{2}), F suits M in a compositional manner — this is necessary for a reanalysis to take place (see \citealt{Eckardt2006}, \citealt{Bar-asher-siegal2020}, \citealt{Aloni2021}). The compositionality principle can be summarised as follows. A complex linguistic expression E such as \textit{The boy is sleeping} is compositional if the meaning of E depends on, and only on, (i) the (morpho)syntactic structure of E and (ii) the meanings of E’s simple linguistic units \citep{Frege1997}. In other words, the meaning of E depends on the meaning of the definite article \textit{the} and the noun \textit{boy}, which combined together syntactically provide a more complex meaning of ‘a particular boy mentioned previously in the discourse’, the meaning of the verbal aspect morphology expressed in \textit{is\,+\,\textup{Verb}\,+\,ing} and the meaning of the verbal lemma \textit{sleep}, which combined together provide the meaning of ‘is sleeping’ and, finally, the combined meaning of the complex expressions has the sentence meaning ‘the boy is sleeping’. The general idea of diachronic formal semantics is that meaning variation and change take place in a transparent way and that every diachronic stage or step in the reanalysis process involves compositionality.


Other linguists take a different perspective assuming that meaning change does not need to be fully compositional or transparent  (\citealt{CroftCruse2004}, \citealt{DeSmetFischer2017}, among many others).

This book has the general goal of testing the compositionality and the transparency principle in formal diachronic semantics on meaning change and variation of the modal indefinite \textit{cualquiera} ‘any’ in Argentinian and European Spanish in synchrony and diachrony. More precisely, my main goal is to investigate the question as to how the modal meaning of the indefinite \textit{cualquiera} ‘any’ changed into evaluative meanings ‘common/unremarkable’ in European and Argentinian Spanish and into the evaluative meaning ‘nonsense’ in Argentinian Spanish, and how this change fits models of linguistic reanalysis in the frameworks of compositional analysis of meaning change in formal diachronic semantics. While acknowledging functionalist accounts of language change, this book places greater emphasis on testing principles of compositionality and formal transparency, focusing on how structural meaning components interact and reconfigure over time.

\section{Structure of the book} \label{sec:1.1}

In \sectref{sec:1.2} and \ref{sec:1.3}, I provide an introduction into the meanings of \textit{cualquiera}, which will define the scope of my study. In \sectref{sec:1.4}, I give a short overview of the current state of the art on meaning change of \textit{cualquiera}, and in \sectref{sec:1.5}, I present very briefly the general methods that will be used to approach my research questions.

  \chapref{sec:2} offers a detailed review of the data and analyses discussed in the literature concerning \textit{cualquiera} in Modern European and Latin American Spanish in synchrony and diachrony. We will see that most synchronic and diachronic approaches do not investigate how the new meanings of \textit{cualquiera} are related to the modal meaning ‘any’ and what might have triggered this meaning change.

  In \chapref{sec:3}, I will present new data and a new analysis of special readings of \textit{cualquiera} in modern European Spanish by extracting new corpus data. %and testing the corpus analysis by speakers judgements from questionnaires and analysing them.

  In \chapref{sec:4}, I will focus on \textit{cualquiera} in Argentinian Spanish in synchrony and diachrony and will suggest a synchronic and diachronic analysis for \textit{cualquiera} in this variety.

  \chapref{sec:5} presents my research on the diachrony of \textit{cualquiera} in European Spanish. In this chapter, I will focus on the distribution of \textit{cualquiera} in various linguistic contexts or variables such as verbal position and verbal properties of the sentence in which \textit{cualquiera} occurs across time. I will raise the question of whether these variables play any role in the meaning change of \textit{cualquiera} across time using a quantitative/statistical analysis. In this chapter, I will provide a formal semantic analysis of different steps that led to meaning change of \textit{cualquiera} and its diachronic development.

  Finally, \chapref{sec:6} summarises the most important findings of this book and their implications for theoretical research on meaning change.

\section{Modal reading of Spanish \textit{cualquiera}}\label{sec:1.2}

The main interest here is the indefinite \textit{cualquiera}, which can roughly be translated with ‘any’ and is hence usually described as a free choice item (FCI) in (certain varieties of) Spanish (see \citealt{Brucart1999}, \citealt{Sanchezlopez1999}, \cite{Alonso-OvalleMenéndez-Benito2011}/\citeyear*{AlonsoMenendez2018}, \citealt{Rivero2011}, \citealt{Eberenz2000}). To introduce this item and before I turn to its semantics (which is the main focus of this book), let us have a brief look at its syntax and morphology:


\ea \label{ex:key:1}
\ea \label{ex:key:1a}
    \gll \textit{cualquier} libro\\  
    \textsc{cualquiera} book\\ %\footnote{{Since} {\textit{cualquiera}} {has different interpretations, as will be shown in this chapter and elaborated on throughout the book, I do not translate it in the glosses.}}
    \glt ‘any book’  

\ex \label{ex:key:1b}
    \gll una camisa \textit{cualquiera}\\
     a shirt \textsc{cualquiera}\\
    \glt ‘any shirt’ (\citealt{AloniEtAl2011}: 2 / DLE, s.v. \textit{cualquiera})
\z
\z


\ea \label{ex:key:2}
    \gll Esto lo puede hacer \textit{cualquiera}.\\
     this it can do \textsc{cualquiera}\\
    \glt ‘Anybody can do this.’ (@ 1\footnote{{The symbol @ indicates that the source of this example is from an internet source. The number represents the index under which the example can be found in the Appendix, where all sources are listed.}})
\z


In \REF{ex:key:1}, \textit{cualquiera} is an invariable indefinite element, which can occur pre- or postnominally. In its prenominal position \REF{ex:key:1a}, it usually drops the \textit{{}-a} and is not preceded by any article, in contrast to the postnominal position, as in \REF{ex:key:1b}, where it is preceded by an indefinite article and cannot drop the \textit{{}-a}. In \REF{ex:key:2}, it is an indefinite pronoun. In this function, the final \textit{{}-a} is always preserved. It does not vary for number and gender, its morphological plural form (\textit{cualesquiera}) being almost obsolete today (\citealt{DeBruyne2011}: 236).

FCIs are indefinites that come along with a free choice implicature or inference (i.e. each individual is a permissible option) (see \citealt{Arregui2006}, \citealt{AloniPort2013}, \citealt{Aloni2021}, \citealt{Alonso-OvalleMenéndez-Benito2011,am2018}, among many others):

\ea \label{ex:key:3}
    \gll \textit{Puedes} \textit{traer} \textit{cualquier} \textit{libro.}\\
    can bring \textsc{cualquier} book\\
    \glt ‘You can bring any book.’\\
    Conventional meaning: ‘You can bring me a book and each book is a possible option.’
    (\citealt{AloniEtAl2011}: 2)
\z

\textit{Cualquier(a)} is described in the literature as a lexically encoded free choice indefinite because its free choice inference seems to be integrated into its semantic content in sentences like \REF{ex:key:3} (see \citealt{AloniPort2013}, \citealt{Aloni2021}, \citealt{Alonso-OvalleMenéndez-Benito2011,am2018}), whereas the free choice inference is possible but not necessarily available with plain indefinites such as \textit{uno} in \REF{ex:key:4} (see \citealt{Aloni2021}):

\ea \label{ex:key:4}
    \gll \textit{Puedes} \textit{traer} \textit{un} \textit{libro.}\\
    can bring a book\\
    \glt ‘You can bring a book.’\\

    a.  Conventional meaning: ‘You can bring a book.’\\
    b.  Free choice inference: ‘Each book is a possible option.’\\
      (\citealt{AloniEtAl2011}: 2)
\z

FCIs are known to appear in modal contexts such as with modal verbs — followed by \textit{que} ‘that’ — to introduce free relative clauses with a subjunctive verb, as in \REF{ex:key:5} (see \citealt{Quer2000}, \citealt{Giannakidou2001}, among many others). This also applies to \textit{cualquiera}:

\ea \label{ex:key:5}
   \gll \textit{Cualquiera} que sea tu caso, vamos a intentar que en este artículo encuentres la respuesta.\\
    \textsc{cualquiera} that be your case go.\textsc{2PL} to try that in this article find.\textsc{2SG} the answer\\
    \glt ‘Whatever your case is, we will try to make sure you find the answer in this article.’

    FC: ‘Each case is a possible option.’

      (@ 2)
\z

Importantly, it has been noted in the literature that \textit{cualquiera} as well as comparable items in other Romance languages can have a series of different other interpretations that diverge from the free choice implicature (cf. \citealt{Alonso-OvalleMenéndez-Benito2011}: 29, \citealt{Rivero2011} for Spanish; \citealt{Fălăuş2015} for \textit{oarecare} in Romanian; and \citealt{Chierchia2006}: 543 for \textit{qualunque} in Italian). As will be shown, this is also the case in Spanish.

  In what follows, I present interpretations other than the free choice interpretation triggered in modal contexts, such as in \REF{ex:key:4} or \REF{ex:key:5}.

\section{Special readings of Spanish \textit{cualquiera}} \label{sec:1.3}

First, the indefinite \textit{cualquiera} can have a random choice interpretation (RCI) in postnominal position with volitional verbs with perfective aspect (\ref{ex:key:6}) (cf. \citealt{ChoiRomero2008}, \citealt{Alonso-OvalleMenéndez-Benito2011}: 29 and \citealt{Rivero2011} for Spanish; \citealt{Fălăuş2015} for \textit{oarecare} in Romanian; \citealt{Chierchia2006}: 543, \citealt{Kellert2021a} for \textit{qualunque/qualsiasi} in Italian; and \citealt{Vlachou2007,Vlachou2012} for \textit{quelconque} in French):

\ea \label{ex:key:6}
    \gll Juan cogió  una  carta  \textit{cualquiera}.\\
    Juan grabbed  a  card  \textsc{cualquiera}\\
    \glt ‘Juan picked a random card.’
    
    Random choice interpretation (RCI): ‘Juan picked a card, and he chose it indiscriminately — he could have picked any other.’
    
    (\citealt{am2018}: 13)
\z

RCI is also known under the term ‘indiscriminative function’, i.e. it conveys the agent’s indifference with respect to his/her choice from a given set of alternatives (for Spanish, see \citealt{am2015}: 14, \citealt{Horn2000} for English; \citealt{Vlachou2007,Vlachou2012} for French). The indiscriminative function has been also observed for \textit{wh}{}-\textit{ever} relatives in English; for further illustration of this function, see the sentence in \REF{ex:key:7}, where \textit{whatever} indicates that the agent’s (in this case the speaker’s) choice of tool was indiscriminate (\citealt{Fintel2000}):

\ea \label{ex:key:7}
    I grabbed whatever tool was handy.\\
      (\citealt{Fintel2000}: 32)
\z

Second, postnominal \textit{cualquiera} can also have an evaluative interpretation (EVAL). It describes the speaker’s (negative) attitude towards some entity, for example that he/she finds Juan ordinary\footnote{{Note that EVAL is a descriptive term that refers to different attitude predicates that have different flavours or readings such as ‘ordinary’, ‘common’, or ‘unremarkable’. These different flavours highly depend on the context (see \sectref{sec:3}).}} (for Spanish, see \citealt{Alonso-OvalleMenéndez-Benito2011}: 29 and  \citealt{Simik2008} for Czech)

\ea \label{ex:key:8}
    \gll Juan es un hombre \textit{cualquiera}. (Spanish)\\
    Juan is a man \textsc{cualquiera}\\
    \glt ‘Juan is an ordinary man.’ (EVAL)\\
    (@ 3)
\z

Third, postnominal \textit{cualquiera} can also have a depreciatory meaning of ‘worthless’ or refer to low quality (see \citealt{Rivero2011}, \citealt{Alonso-OvalleMenéndez-Benito2011} for Spanish; \citealt{Fălăuş2015} for Romanian \textit{oarecare}; \citealt{Vlachou2007,Vlachou2012} on French and Greek FCIs). I consider this meaning to be a subset of the evaluative interpretation (EVAL):

\ea \label{ex:key:9}
    \gll Es un hombre \textit{cualquiera.} (Spanish)\\
    is a man \textsc{cualquiera}\\
    \glt ‘He is a man without great qualities.’ (EVAL)\\
    (@ 4)
\z

The evaluative reading is particularly interesting, because it leads to further processes of recategorisation and lexicalisation. For example, the French FCI \textit{quelconque} ‘whichever’ shows a shift from a free choice indefinite to a gradable adjective:

\ea \label{ex:key:10}
    \gll Elle m’a paru assez \textit{quelconque}. (French)\\
    she {me has} looked quite \textsc{quelconque}\\
    \glt ‘She looked quite common to me.’\\
    \citep[1530]{Vlachou2012}
\z

Similarly, the French FCI \textit{n’importe quoi} (‘whatever’, lit. ‘does not matter what’) can also be used to express EVAL. In contrast to the adjectival use of \textit{quelconque} in \REF{ex:key:10}, \textit{n’importe} \textit{quoi} is a noun that can be modified by an adjective and has the meaning ‘nonsense/bullshit’:

\ea \label{ex:key:11}
    \gll Il raconte du grand \textit{n’importe quoi}. (French)\\
    he tells of-the big \textsc{n’importe quoi}\\
    \glt ‘He talks real bullshit.’\\
    (@ 138)
\z

In the course of this book, I will show similar effects for Argentinian Spanish (ArgSp), where \textit{cualquiera} has lexicalised the meaning EVAL and can nowadays be used as a gradable adjective modifiable by degree adverbs or the degree modifier \textit{muy} or \textit{re} (\citealt{Rizzosalierno2013}: 27):

\ea \label{ex:key:12}
    \gll La peli es {(re/muy)} \textit{cualquiera}. (ArgSp)\\
    the movie is { very} \textsc{cualquiera}\\
    \glt ‘The movie is (very) unremarkable/bad/not of great interest.’\footnote{The degree modifier \textit{re} is used in Argentinian Spanish to reinforce the meaning of linguistic expressions. When \textit{re} is combined with gradable adjectives, it has the meaning of ‘very’ (see \citealt{Rizzosalierno2013} and \chapref{sec:4} for further discussion).}\\
    (@ 5 / @ 6)
\z

\textit{Cualquiera} in Argentinian Spanish can also be used as a noun under transitive verbs, similar to \textit{n’importe quoi} in \REF{ex:key:11} with the meaning ‘nonsense/bullshit’:

\ea \label{ex:key:13}
    \gll Me dijiste \textit{cualquiera}. (ArgSp)\\
    me said \textsc{cualquiera}\\
    \glt ‘You told me bullshit/nonsense.’\\
    (@ 7)
\z

In this book, I distinguish between \textit{cualquiera} with the interpretation of ‘unremarkable, bad, not of great interest’, which is mostly found as a predicative complement as in \REF{ex:key:12} (henceforth: EVAL1) and \textit{cualquiera} embedded under transitive verbs with the meaning ‘bullshit/nonsense’ as in \REF{ex:key:13} (henceforth: EVAL2).

We find the same lexicalisation pattern in Mexican Spanish for the lexical item \textit{equis} (originally meaning ‘the letter x’), which can be used as an indefinite with a free choice interpretation ‘any’ in conditional clauses as in \REF{ex:key:14} as shown in \citet[137]{Kellert2020semantic}, but can also function as a gradable adjective that expresses EVAL1 as shown in \REF{ex:key:15} (see \citealt{Kellert2020semantic}: 139):

\ea \label{ex:key:14}
    \gll Si por \textit{equis} razón no vas a venir, avísame  \hfill (MexSp)\\
if for X reason \textsc{neg} go.\textsc{2SG} to come let.me.know\\
    \glt ‘If you cannot come for whatever reason, inform me.’

    \citep[137]{Kellert2020semantic}
\z

\ea \label{ex:key:15}
    \gll un hombre bastante/ demasiado/ muy equis  \hfill (MexSp)\\
a man quite/ too/ very X\\
    \glt ‘a very ordinary man (with no special qualities)’

    \citep[138]{Kellert2020semantic}
\z

Finally, \textit{cualquiera} can be interpreted as Negative Polarity Item (NPI) in European Spanish, i.e. as an existential indefinite in the scope of negation with the interpretation of ‘not … any’ (see \citealt{DeBruyne2011}, §497, \citealt{Arregui2006}, \citealt{Chierchia2006}, among others):

\ea \label{ex:key:16}
    \gll Estos últimos, además de no comer carne y pescado, no comen \textit{cualquier} cosa que sea de origen animal, como la miel, la leche y los huevos.\\
these last.ones besides of \textsc{neg} eating meat and fish \textsc{neg} eat \textsc{cualquier} thing that be of origin animal like the honey the milk and the eggs\\
    \glt ‘The latter ones, apart from not eating meat and fish, don’t eat any animal products such as honey, milk, and eggs.’

    (@ 8)
\z

However, the NPI interpretation of \textit{cualquiera} as ‘anything’ is restricted to relative clause modification, because without relative clause modification, \textit{cualquiera} has the meaning ‘everything’ or a ‘not just any’ interpretation under negation (see \citealt{Arregui2006}, among others; for the general phenomonen, see \citealt{Chierchia2006}):

\ea \label{ex:key:17}
    \gll No comen \textit{cualquier} cosa.\\
    \textsc{neg} eat \textsc{cualquier} thing\\
    \glt ‘They don’t eat everything.’ or ‘They don’t eat just anything. They eat something special.’\\
    \# ‘They don’t eat anything.’ \\
    (@ 9)
\z

Given all the meanings in the examples from \REF{ex:key:1} to \REF{ex:key:17}, the question arises as to whether these meanings reflect lexical ambiguity or polysemy. In this case, each reading of \textit{cualquiera} has a different denotation (as with the word \textit{paper} in English, which can denote a newspaper, an academic publication, or the physical material itself) (see Hypothesis 1).
In the other case, \textit{cualquiera} has only one denotation and the different meanings described above are the result of some pragmatic inference on top of the semantic denotation (see Hypothesis 2):


\ea \label{ex:key:18}
    Hypothesis 1: \textit{Cualquiera} has different denotations (a case of polysemy).

    Hypothesis 2: \textit{Cualquiera} has one denotation and several readings follow from this denotation.
\z

In this book, I will argue that the latter case is true for European Spanish and that \textit{cualquiera} has only one denotation there, namely the denotation of an existential quantifier with the meaning ‘some (or other)’ (see \citealt{Fălăuş2014,Aloni2021}, among others for analysing FCIs as existential quantifiers), whereas it is lexically ambiguous in Argentinian Spanish (‘some or other’ and ‘bad’).

Despite some progress in the research on free choice indefinites, there is still no agreement about the semantic relation between RCI and EVAL1 (\citealt{am2015}: 4), that is, on whether these interpretations or at least some subset of them have a common semantic source (i.e. semantic denotation) and, if so, how one interpretation is related to another. \citet{Alonso-OvalleMenéndez-Benito2011} assume as a working hypothesis that \textit{cualquiera} is ambiguous between the random choice interpretation (RCI) and the evaluative interpretation as ‘unremarkable’ (EVAL1) in Modern Spanish. However, as they mention themselves, this working hypothesis should be studied in detail (\citealt{Alonso-OvalleMenéndez-Benito2011}: 41). In this work, I will test this hypothesis. Moreover, I will also study \textit{cualquiera} in Argentinian Spanish, which allows EVAL2, in contrast to European Spanish, as will be shown in \sectref{sec:4}.

\section{Open issues in previous research and goals of the present book} \label{sec:1.4}

Most synchronic accounts of \textit{cualquiera} focus on one of its particular interpretations such as the random choice interpretation (see \citealt{ChoiRomero2008}, \citealt{Alonso-OvalleMenéndez-Benito2011}) or the free choice interpretation (see \citealt{Menéndez-Benito2010}, \citealt{AloniPort2013}, \citealt{Fălăuş2015}, among others). There is no systematic investigation of the distribution and restrictions of the evaluative meaning of \textit{cualquiera}. Moreover, previous accounts do not try to unify the meanings that we are concerned with in this investigation or to derive one meaning from the other. It thus remains to be seen whether the evaluative meaning can be derived from the free choice function of indefinites.

The main aim of this book is to fill this gap by offering a new analysis that can unify different meanings of \textit{cualquiera} that look unrelated in synchrony, namely the FCI and RCI meanings ‘any’ and ‘random’ and the EVAL meanings ‘ordinary, bad, common, nonsense’. In order to achieve this goal, it will be crucial to incorporate a diachronic perspective on meaning variation. In fact, as I will demonstrate, it turns out that at some diachronic stage, EVAL1 was derived from FCI through a certain syntactic and semantic reanalysis of FCIs. By looking at the history of the meanings and functions of \textit{cualquiera} we will be able to identify how one particular meaning evolved from the other. In sum, a diachronic perspective can help us to reconstruct the relation between different meanings in a compositional way.

This perspective of combining synchrony and diachrony has not been discussed so far in the literature on the semantic evolution of \textit{cualquiera}. The few existing diachronic studies concentrate on the question of how the pronoun \textit{qual} (\textit{cual} in Modern Spanish) and the verb form \textit{quiera} (want.\textsc{3sg.pres.sbjv}) became grammaticalised into one complex compound (see \citealt{Palomo1934}, \citealt{Lombard1938}–39/47–48, \citealt{Rivero1988}, \citealt{Company-CompanyPozas-Loyo2009}, \citealt{Mackenzie2019}, among others). These accounts, however, do not discuss the question of what consequences this grammaticalisation process has for the interpretational component of \textit{cualquiera}. Furthermore, it is still unclear when and how the different meanings associated with this form have emerged.

Moreover, there is no systematic quantitative analysis of diachronic data on \textit{cualquiera} that includes a statistical significance test providing empirical support for a change. While \citet{cp2009} and \citet{AloniPort2013} have presented a quantitative description of \textit{cualquiera}, they did not use any statistical test to verify their hypotheses. Also, the classification of the data in the literature is not satisfactory. \citet{AloniPort2013} examine a number of syntactic and semantic properties of \textit{cualquiera}, but they do not attempt to establish a correlation between the different factors considered. From their results we do not learn whether different functions correlate with different syntactic positions of \textit{cualquiera} (e.g. +/– prenominal position, +/– preverbal, +/– appearance in episodic context or copular constructions).

In view of the shortcomings of previous studies, it seems relevant to use a more detailed syntactic classification for a corpus analysis of synchronic and diachronic data on \textit{cualquiera}. In particular, I will look at the morphosyntactic and semantic properties of the verb forms that co-occur in the sentences in which \textit{cualquiera} appears and will describe them with respect to their modality, lexical specification, tense, and aspect. This will allow us to observe correlations between the different meanings of \textit{cualquiera} and the aforementioned verbal properties. As for the diachrony of \textit{cualquiera}, I will investigate the first occurrences of its different meanings in order to draw conclusions about the order in which these meanings appeared over time.

In light of these empirical and theoretical gaps, this book aims not only to unify the different meanings of \textit{cualquiera} through a diachronic compositional analysis but also to make a methodological contribution by grounding this analysis in a combination of annotated corpus data, statistical evaluation, and informal native speaker consultation. While most previous work has focused on isolated interpretations of \textit{cualquiera}, this study systematically traces the development of its modal and evaluative readings across time and regional varieties. The approach taken here relies primarily on principles of compositionality and transparency in formal semantics rather than functionalist or discourse-based explanations. Through this lens, the book explores how shifts in meaning can be modelled as semantically and syntactically motivated reanalyses. In doing so, it presents the first comprehensive monograph on \textit{cualquiera} that integrates corpus-based evidence, diachronic distribution, and theoretical modelling to account for both synchronic variation and historical change.

\section{Data and methodology} \label{sec:1.5}

In order to describe the evaluative properties of \textit{cualquiera} in Modern European Spanish, I have extracted data with a noun + \textit{cualquiera} (henceforth (\textsc{un)} N \textit{cualquiera}) from the \textit{Corpus del Español web dialects} (CdE\_web), which is annotated for different regional varieties of Spanish. The data on the postnominal as well as on the pronoun \textit{cualquiera} in Argentinian Spanish and in other Latin American varieties including European Spanish was extracted from the corpus \textit{Corpus del Español}. 

In order to respond to my diachronic research question, I have chosen 150 randomised occurences per century from the 13th to the 20th centuries from the \textit{Corpus del Español\_hist-gen} (henceforth CdE\_hist-gen) and have annotated them according to 30 syntactic and semantic properties. As some of the texts contained in CdE\_hist-gen are taken from editions that have low philological reliability, as reported in \citet{OctaviodeToledo2017}, I selected only those texts that these authors classify as suitable for linguistic analysis. {In addition, I used particular search algorithms, which will be specified in the respective parts of the description of the corpus analysis. These algorithms were designed to target approximations of syntactic and semantic properties of my small sample in the whole corpus CdE\_hist-gen. The goal of this task was to see whether it is possible to identify similar linguistic patterns of diachronic change in the whole corpus without taking into account manual annotation and the distinction between texts regarding their philological reliability.} For the diachronic research in Argentinian Spanish, I used the diachronic corpus CORDIAM, \textit{Corpus Diacrónico y Diatópico del Español}, which contains a total of 10,889,176 words from 19 Latin-American countries from 1494 to 1905, with documents from archives, journals, and literature. It is annotated for geographical origins. The corpus CORDIAM does not contain many instances of \textit{cualquiera} in Argentinian Spanish in comparison to CdE\_hist-gen (around 150 in total). I therefore also used a larger diachronic corpus, CORDE (\textit{Corpus Diacrónico del Español}), which is a diachronic corpus with written texts spanning from Old Spanish to 1974 with annotation of geographic origin (country), and which contains 544 instances of \textit{cualquiera} in Argentinian Spanish. I analysed all diachronic data from Argentinian Spanish from CORDIAM and CORDE as the data collection was relatively small.

The data I analysed for the quantitative analysis can be found on GitHub: \url{https://github.com/olga-kel/Historical-Linguistics}.

The quantitative analysis was based on statistical methods such as the chi-square test and the linear regression model (see \citealt{Baayen2008}). I have used the R software programme for the statistical analysis of the data. The generation of the plots and histograms was carried out in Microsoft Excel.

{Corpora are the main source of data in this book. In addition to the corpus data, I informally consulted several native speakers of European and Argentinian Spanish. These consultations were not part of a systematic experimental survey but served to validate the corpus-based analysis of the use and interpretation of \textit{cualquiera} through selective, informal interviews. These interviews were conducted with linguistics students from Spain, as well as with members of the Facebook group \textit{Lingüística Argentina}, where both students and instructors with a background in linguistics responded to my questions. The selection of informants with linguistic training is based on the assumption that such individuals tend to be more reliable language informants, as they are often better able to articulate and reflect on language use. The informal judgments were used to assess whether speakers’ intuitions align with the corpus-based findings, which remain the central foundation of this study. Future experimental work will be necessary to more systematically describe the distribution of \textit{cualquiera} in both European and Argentinian Spanish.}   

For the semantic analysis of synchronic and diachronic data, I used formal semantic frameworks of FCIs (\citealt{Chierchia2006}, \citealt{Aloni2021}, \citealt{AloniPort2013}, \citealt{Alonso-OvalleMenéndez-Benito2011,am2018}, among others) and models of semantic reanalysis within formal semantics (\citealt{Eckardt2006}, \citealt{AloniPort2013}, among others).

